\documentclass[10pt, letterpaper]{article}

% Packages:
\usepackage[
    ignoreheadfoot, % set margins without considering header and footer
    top=2 cm, % seperation between body and page edge from the top
    bottom=2 cm, % seperation between body and page edge from the bottom
    left=2 cm, % seperation between body and page edge from the left
    right=2 cm, % seperation between body and page edge from the right
    footskip=1.0 cm, % seperation between body and footer
    % showframe % for debugging 
]{geometry} % for adjusting page geometry
\usepackage{titlesec} % for customizing section titles
\usepackage{tabularx} % for making tables with fixed width columns
\usepackage{array} % tabularx requires this
\usepackage[dvipsnames]{xcolor} % for coloring text
\definecolor{primaryColor}{RGB}{0, 79, 144} % define primary color
\usepackage{enumitem} % for customizing lists
\usepackage{fontawesome5} % for using icons
\usepackage{amsmath} % for math
\usepackage[
    pdftitle={John Doe's CV},
    pdfauthor={Malthe Sporring},
    pdfcreator={LaTeX with RenderCV},
    colorlinks=true,
    urlcolor=primaryColor
]{hyperref} % for links, metadata and bookmarks
\usepackage[pscoord]{eso-pic} % for floating text on the page
\usepackage{calc} % for calculating lengths
\usepackage{bookmark} % for bookmarks
\usepackage{lastpage} % for getting the total number of pages
\usepackage{changepage} % for one column entries (adjustwidth environment)
\usepackage{paracol} % for two and three column entries
\usepackage{ifthen} % for conditional statements
\usepackage{needspace} % for avoiding page brake right after the section title
\usepackage{iftex} % check if engine is pdflatex, xetex or luatex

% Ensure that generate pdf is machine readable/ATS parsable:
\ifPDFTeX
    \input{glyphtounicode}
    \pdfgentounicode=1
    % \usepackage[T1]{fontenc} % this breaks sb2nov
    \usepackage[utf8]{inputenc}
    \usepackage{lmodern}
\fi



% Some settings:
\AtBeginEnvironment{adjustwidth}{\partopsep0pt} % remove space before adjustwidth environment
\pagestyle{empty} % no header or footer
\setcounter{secnumdepth}{0} % no section numbering
\setlength{\parindent}{0pt} % no indentation
\setlength{\topskip}{0pt} % no top skip
\setlength{\columnsep}{0cm} % set column seperation
\makeatletter
\let\ps@customFooterStyle\ps@plain % Copy the plain style to customFooterStyle
\patchcmd{\ps@customFooterStyle}{\thepage}{
    \color{gray}\textit{\small Malthe Sporring - Page \thepage{} of \pageref*{LastPage}}
}{}{} % replace number by desired string
\makeatother
\pagestyle{customFooterStyle}

\titleformat{\section}{\needspace{4\baselineskip}\bfseries\large}{}{0pt}{}[\vspace{1pt}\titlerule]

\titlespacing{\section}{
    % left space:
    -1pt
}{
    % top space:
    0.3 cm
}{
    % bottom space:
    0.2 cm
} % section title spacing

\renewcommand\labelitemi{$\circ$} % custom bullet points
\newenvironment{highlights}{
    \begin{itemize}[
        topsep=0.10 cm,
        parsep=0.10 cm,
        partopsep=0pt,
        itemsep=0pt,
        leftmargin=0.4 cm + 10pt
    ]
}{
    \end{itemize}
} % new environment for highlights

\newenvironment{highlightsforbulletentries}{
    \begin{itemize}[
        topsep=0.10 cm,
        parsep=0.10 cm,
        partopsep=0pt,
        itemsep=0pt,
        leftmargin=10pt
    ]
}{
    \end{itemize}
} % new environment for highlights for bullet entries


\newenvironment{onecolentry}{
    \begin{adjustwidth}{
        0.2 cm + 0.00001 cm
    }{
        0.2 cm + 0.00001 cm
    }
}{
    \end{adjustwidth}
} % new environment for one column entries

\newenvironment{twocolentry}[2][]{
    \onecolentry
    \def\secondColumn{#2}
    \setcolumnwidth{\fill, 4.5 cm}
    \begin{paracol}{2}
}{
    \switchcolumn \raggedleft \secondColumn
    \end{paracol}
    \endonecolentry
} % new environment for two column entries

\newenvironment{header}{
    \setlength{\topsep}{0pt}\par\kern\topsep\centering\linespread{1.5}
}{
    \par\kern\topsep
} % new environment for the header



% save the original href command in a new command:
\let\hrefWithoutArrow\href

% new command for external links:
\renewcommand{\href}[2]{\hrefWithoutArrow{#1}{\ifthenelse{\equal{#2}{}}{ }{#2 }\raisebox{.15ex}{\footnotesize \faExternalLink*}}}


\begin{document}
    \newcommand{\AND}{\unskip
        \cleaders\copy\ANDbox\hskip\wd\ANDbox
        \ignorespaces
    }
    \newsavebox\ANDbox
    \sbox\ANDbox{}

    %\placelastupdatedtext
    \begin{header}
        \textbf{\fontsize{24 pt}{24 pt}\selectfont Malthe Sporring}

        \vspace{0.3 cm}

        \normalsize
        \mbox{\hrefWithoutArrow{mailto:malthe.sporring@ed.ac.uk}{\color{black}{\footnotesize\faEnvelope[regular]}\hspace*{0.13cm}malthe.sporring@ed.ac.uk}}%
        \kern 1 cm%
        \textbf{Web:}  \mbox{\hrefWithoutArrow{https://malthefogsporring.github.io/}{\color{black}malthefogsporring.github.io}}
    \end{header}

    \vspace{0.3 cm - 0.3 cm}

    \section{Research interests}
    Homotopy theory, ($\infty$-)category theory, algebraic topology, factorisation algebras.

    \section{Education}



        
        \begin{twocolentry}{
            
            
        \textit{Sep 2022 - Aug 2027 (Expected)}}
            \textbf{University of Edinburgh}

            \textit{PhD Geometry and Topology}
        \end{twocolentry}

        \vspace{0.10 cm}
        \begin{onecolentry}
            \begin{highlights}
                \item Working with Prof. Clark Barwick on a project on "$n$-fold monoidal $\infty-$categories."
                \item (Sep - Oct 2025) Received grants from the London Mathematical Society and the Institute for Mathematics and its Applications for a two-month research visit to NTNU Trondheim.
            \end{highlights}
        \end{onecolentry}
        \vspace{0.10 cm} 
        \begin{twocolentry}{
           
            
        \textit{Sep 2017 - May 2022}}
            \textbf{University of Edinburgh}

            \textit{MMath Mathematics}
        \end{twocolentry}

        \vspace{0.10 cm}
        \begin{onecolentry}
            \begin{highlights}
                \item Awarded a 2021 Summer Vacation Scholarship for a project on “Axiomatic Homology and the Borsuk-Ulam Theorem”.
                %\item Awarded a Year 1 Certificate of Merit.
            \end{highlights}
        \end{onecolentry}

    \section{Invited talks}
    \begin{highlights}
    \item \begin{twocolentry}{
        \textit{September 2026}}
            \textbf{Pre-monoidal $\infty$-operads }

            \textit{Universitat Autònoma de Barcelona Topology Seminar}
        \end{twocolentry}
\item \begin{twocolentry}{
        \textit{June 2026}}
            \textbf{The pre-monoidal envelope of an $\infty$-operad} (Poster; available \hrefWithoutArrow{https://malthefogsporring.github.io/files/copenhagen_poster.pdf}{here}.)

            \textit{Young Topologist Meeting, Copenhagen}
        \end{twocolentry}
    \item \begin{twocolentry}{
           
            
        \textit{March 2026}}
            \textbf{Twofold monoidal $\infty$-categories }

            \textit{Warwick University Junior Algebra Seminar}
        \end{twocolentry}
    \item \begin{twocolentry}{
           
            
        \textit{May 2025}}
            \textbf{Twofold monoidal structures}

            \textit{University of Aberdeen Postgraduate Research Seminar}
        \end{twocolentry}
\end{highlights}
    \section{Teaching}

        \begin{twocolentry}{
 
            
        \textit{Jan 2023 - Current}}
            \textbf{Mathematics Tutor}
            
            \textit{University of Edinburgh}
        \end{twocolentry}

        \vspace{0.10 cm}
        \begin{onecolentry}
            \begin{highlights}
                \item (Autumn 2026) Introduction to Mathematics at University
                \item (Spring 2026) Honours Algebra, Honours Algebra - Skills (Mathematical computing).
                \item (Spring 2025) Calculus and its Applications.
                \item (Autumn 2024) Accelerated Proofs and Problem Solving, Axiomatic Set Theory, MathsBase (Drop-in study space).
                \item (Spring 2024) Honours Algebra - Skills (Mathematical computing), Proofs and Problem Solving.
                \item (Autumn 2023) Category Theory, Honours Analysis - Skills (Presentation skills), MathsBase (Drop-in study space).
                \item (Spring 2023) Algebraic Topology, Fundamentals of Pure Mathematics.
            \end{highlights}
        \end{onecolentry}

            \begin{twocolentry}{
        \textit{Fall 2023 and 2025}}
         \vspace{0.2 cm}
            \textbf{Volunteering}
        \end{twocolentry}

        \vspace{0.10 cm}
        \begin{onecolentry}
            \begin{highlights}
                \item Mentored two projects for \hrefWithoutArrow{https://sites.google.com/view/twoples/home}{Twoples}, an online directed reading program for mathematics undergraduates. The first project was on algebraic topology, the second project is on homological algebra.
            \end{highlights}
        \end{onecolentry}
        
    \section{Professional Experience}
    
        \begin{twocolentry}{
        \textit{July - August 2025}}
            \textbf{Online Assessment Intern}
            
            \textit{University of Edinburgh}
        \end{twocolentry}
        \vspace{0.10 cm}
        \begin{onecolentry}
            \begin{highlights}
                \item Took part in a project to turn the HELM mathematics workbooks into STACK quizzes. This included managing interns and performing quality assurance of their work. STACK is an online assessment system for mathematics and STEM.
            \end{highlights}
        \end{onecolentry}
\vspace{0.2 cm}

        \begin{twocolentry}{
        \textit{Aug - Sep 2022}}
            \textbf{Online Assessment Intern}
            
            \textit{University of Glasgow}
        \end{twocolentry}

        \vspace{0.10 cm}
        \begin{onecolentry}
            \begin{highlights}
                \item Responsible for converting a series of diagnostics tests for undergraduate students into the STACK system.
            \end{highlights}
        \end{onecolentry}

        \vspace{0.2 cm}

        \begin{twocolentry}{
        \textit{May - June 2022}}
            \textbf{Online Assessment Intern}
            
            \textit{Heriot-Watt University}
        \end{twocolentry}

        \vspace{0.10 cm}
        \begin{onecolentry}
            \begin{highlights}
                \item Took part in a project to introduce STACK assessment in first year mathematics courses. This included designing new questions, and running a workshop for colleagues.
            \end{highlights}
        \end{onecolentry}


        \vspace{0.2 cm}

        \begin{twocolentry}{
        \textit{Jan - May 2022}}
            \textbf{Freelance Mathematics Tutor}        \end{twocolentry}

        \vspace{0.10 cm}
        \begin{onecolentry}
            \begin{highlights}
                \item Gained experience tutoring students age 16-18 of the International Baccalaureate programme, in both an individual and group setting.
            \end{highlights}
        \end{onecolentry}


        \vspace{0.2 cm}

        \begin{twocolentry}{
        \textit{June - Aug 2020}}
            \textbf{ASID Intern}
            
            \textit{University of Edinburgh}
        \end{twocolentry}

        \vspace{0.10 cm}
        \begin{onecolentry}
            \begin{highlights}
                \item Took part in a project to turn the HELM mathematics workbooks into STACK quizzes. This included conducting quality-assurance on questions written by colleagues.
                \item Gained experience in teaching STACK from organizing and speaking at eight STACK workshops.
                \item Learned HTML and CSS through designing the official website \hrefWithoutArrow{https://stack-assessment.org/}{stack-assessment.org}. 
            \end{highlights}
        \end{onecolentry}

        \vspace{0.2 cm}

        \begin{twocolentry}{
        \textit{June - Aug 2019}}
            \textbf{Online Assessment Community Engagement
Intern}
            
            \textit{University of Edinburgh}
        \end{twocolentry}

        \vspace{0.10 cm}
        \begin{onecolentry}
            \begin{highlights}
                \item Responsible for improving the learning experience of the online mathematics assessment system STACK.
                \item Collated a  \hrefWithoutArrow{https://docs.stack-assessment.org/content/2019-cate-case-studies.pdf}{a collection of cases studies} detailing innovative use of STACK. 
                \item Conducted a survey of STACK users world-wide.
            \end{highlights}
        \end{onecolentry}
    
    
    \section{Technologies}

        \begin{onecolentry}
            Experienced user of Python, Git, Lean and LaTeX.
            
            \vspace{0.2 cm}

            Experienced in online mathematics assessment via STACK
        \end{onecolentry}

    

\end{document}